\documentclass[12pt]{article}
\usepackage[utf8]{inputenc}
\usepackage[T1]{fontenc}
\usepackage{hyperref}
\title{Tillamook Language Resources for K-12 Curriculum Development}
\author{Compiled by Claude Sonnet 5 (claude-sonnet-5) for Andr\'as Kornai}
\date{}
\begin{document}
\maketitle

\section{Introduction}
Tillamook (ISO 639-3 til; Glottocode till1254) is an extinct Salishan language of the Oregon coast. One database lists it as extinct since 1972. There are no native speakers, so a curriculum would have to be a revitalization curriculum built on archival and published documentation. Only a few resources turned up in search, and no classroom materials were among them.

\section{Primary Educational Resources}
No primers, textbooks or school programmes for Tillamook turned up in search. The online talking dictionary below is the closest thing to a teaching tool.

\section{Linguistic and Reference Materials}
\begin{itemize}
\item Tillamook--English Online Talking Dictionary (UBC RelLex), a bilingual, bidirectional word list launched in 2012 with the Tillamook community and the Living Tongues Institute. Its data come from the unpublished Tillamook dictionary of Laurence C. Thompson and M. Terry Thompson. \url{https://wiki.ubc.ca/Documentation:RelLex/Tillamook_-_English_Online_Talking_Dictionary}
\item Reichard's multi-volume comparison of five Salish languages (UBC RelLex), which covers Tillamook. It is a comparative work, not a Tillamook dictionary. \url{https://wiki.ubc.ca/Documentation:RelLex/A_Comparison_of_Five_Salish_Languages_(Volumes_I_-_VI)}
\item Glottolog reference record. \url{https://glottolog.org/resource/reference/id/31549}
\item ASJP word-list entry. \url{https://asjp.clld.org/languages/TILLAMOOK}
\end{itemize}

\section{Religious and Community Resources}
Background on the Tillamook people. This is not language material. \url{https://en.wikipedia.org/wiki/Tillamook_people}

\section{Digital and Audio Resources}
The talking dictionary above is the only audio-enabled resource found. The Endangered Languages Project also lists a Tillamook sample. \url{https://www.endangeredlanguages.com/lang/5396/samples/11149?hl=de}

\section{Curriculum Development Context}
Documentation is thin, and the main dictionary is based on unpublished materials. A curriculum would need to work with the Tillamook community and archives. It could start with vocabulary, place names and cultural content from the talking dictionary. A library catalogue record seen in search may point to further print sources. \url{https://opac.nwic.edu/eg/opac/record/36971?locg=1}
\end{document}